% DeepPass68: general d-dimensional integral identities for loop calculations.

$d$ 维径向积分需要单位球面的面积。由定义
\begin{equation}
\Omega_d\equiv\int\dd\Omega_{d-1},
\label{eq:d-sphere-solid-angle-definition-dp68}
\end{equation}
可得
\begin{equation}
\boxed{
\Omega_d=\frac{2\pi^{d/2}}{\Gamma(d/2)}.}
\label{eq:d-sphere-solid-angle-dp21}
\end{equation}

把 Schwinger 参数式用于 Euclidean 分母，先写成
\begin{equation}
\frac1{(\ell_E^2+\Delta)^\nu}
=\frac1{\Gamma(\nu)}
\int_0^\infty\dd s\,s^{\nu-1}\ee^{-s(\ell_E^2+\Delta)}.
\label{eq:euclidean-master-schwinger-dp68}
\end{equation}
在 $\Re\Delta>0$、$\Re\nu>d/2$ 的收敛域内先完成 $\ell_E$ 与 $s$ 积分，得到
\begin{equation}
\boxed{
\int\frac{\dd^d \ell_E}{(2\pi)^d}
\frac1{(\ell_E^2+\Delta)^\nu}
=\frac1{(4\pi)^{d/2}}
\frac{\Gamma(\nu-d/2)}{\Gamma(\nu)}
\Delta^{d/2-\nu}.}
\label{eq:euclidean-master-integral-dp21}
\end{equation}
右端随后定义对 $d$、$\nu$ 与 $\Delta$ 的解析延拓。

四阶张量分子需要比\eqnrefs{eq:tensor-reduction-proof-r9} 多一个对称结构。完全对称且各向同性的秩四张量只能写成
\begin{equation}
\int\dd^d \ell\,\ell^\mu\ell^\nu\ell^\rho\ell^\sigma f(\ell^2)
=C_4\bigl(
\eta^{\mu\nu}\eta^{\rho\sigma}
+\eta^{\mu\rho}\eta^{\nu\sigma}
+\eta^{\mu\sigma}\eta^{\nu\rho}\bigr).
\label{eq:rank4-isotropic-ansatz-dp68}
\end{equation}
双重收缩固定系数后得到
\begin{equation}
\boxed{
\int\dd^d \ell\,
\ell^\mu\ell^\nu\ell^\rho\ell^\sigma f(\ell^2)
=\frac{\eta^{\mu\nu}\eta^{\rho\sigma}
+\eta^{\mu\rho}\eta^{\nu\sigma}
+\eta^{\mu\sigma}\eta^{\nu\rho}}
{d(d+2)}
\int\dd^d \ell\,(\ell^2)^2f(\ell^2).}
\label{eq:rank4-tensor-reduction-dp21}
\end{equation}

在 $d=4-2\epsilon$ 附近，圈积分中经常出现同一个无量纲组合
\begin{equation}
\mathcal U_\epsilon(\Delta)
\equiv\Gamma(\epsilon)(4\pi)^\epsilon
\left(\frac{\Delta}{\mu^2}\right)^{-\epsilon},
\label{eq:dimreg-universal-factor-dp68}
\end{equation}
其中 $\mu$ 是保持重整化参数具有固定工程量纲所引入的尺度。小 $\epsilon$ 展开为
\begin{equation}
\boxed{
\mathcal U_\epsilon(\Delta)
=\frac1\epsilon-\gamma_E+\ln4\pi
-\ln\frac{\Delta}{\mu^2}+\order(\epsilon).}
\label{eq:dimreg-universal-expansion-dp68}
\end{equation}
因此 $\overline{\rm MS}$ 中反复出现的 $1/\epsilon-\gamma_E+\ln4\pi$ 只是这个通用展开的方案常数部分。
